\documentclass[10pt,a4paper]{notes-template}

\usepackage[
  backend=biber,
  hyperref=true,
  backref=true,
  giveninits=true,
  style=alphabetic,
  sortcites=false,
  maxbibnames=99,
  maxcitenames=2,
  sorting=ydnt,
  defernumbers=true,
]{biblatex}

\addbibresource{biblio.bib} % Bibliography file
\AtEveryBibitem{\clearfield{isbn}}
\AtEveryBibitem{\clearfield{issn}}
\AtEveryBibitem{\clearfield{volume}}
\AtEveryBibitem{\clearlist{location}}
\AtEveryBibitem{\clearfield{series}}

\usepackage{tikz}
\usepackage{tikz-cd}
\tikzset{>={Straight Barb[scale=0.8]}}
\tikzcdset{
arrow style=tikz
}

\usepackage{macros}

\begin{document}

  \title{Interpolation for the simply-typed λ-calculus}

  \author{Meven \textsc{Lennon-Bertrand} \and Alexis \textsc{Saurin}}

\maketitle

\begin{abstract}
  Rewriting \textcite{Cubric1994} in a modern type-theoretic language.
\end{abstract}

\section{Context}

\textcite{Cubric1994} proves interpolation for the internal language of biCCC,
that is, \STLC with at least pair, unit, sum, and empty types, all with η-laws
(including for positive connectives!), plus arbitrary base types, constants,
and equations.

He really needs commuting conversions for \(+\) and \(\bot\), apparently because
of some subformula property. I wonder how this looks like in our setting…

\section{Definitions}

\subsection{λ-calculus}

\paragraph{Definition}
We work with \STLC, which we take to always have unit, product and function types
(\ie this is the internal language of CCC).
% We write \STLCp for the addition
% of an empty type and pair type to the above. If we throw η into the mix, this is
% the internal language of biCCC. But maybe we want to negotiate on these, which
% don't fare very well in the dependently-typed setting…
%
We moreover work with respect to a theory \(\Theta\).

\begin{definition}{Language}{lang}
  A \emph{language} \(\Lambda\) consists of:
  \begin{itemize}
    \item a set of base types \(\mathcal{B} \in \Set\);
    \item a family of typed constants indexed by types
      \(\mathcal{C} \in \Ty_{\mathcal{B}} \to \Set\);    
  \end{itemize}
  In the above, \(\Ty_{\mathcal{B}}\) is the set of types generated by
  \(\mathcal{B}\) which will be defined in \cref{fig:stlc}.
  If \(\Lambda = (\mathcal{B},\mathcal{C})\) we write
  \(\Ty_{\Lambda}\) for \(\Ty_{\mathcal{B}}\), and
  we write \(\emptyset\) for the empty language, \ie that where the sets
  of base types and constants are empty.
\end{definition}

Technically, we should interleave
the definition of types generated by a set of base types
with that of a language, but that's annoying
(and becomes impossible anyway once you move to dependent type theory, see
\eg \textcite{Bezem2021} for the issue and an elegant solution).

\begin{definition}{Theory}{sig}
  A \emph{theory} \(\Theta\) consists of:
  \begin{itemize}
    \item a language \(\Lambda\);
    \item a set of equations \(\mathcal{E} \ty
      (A : \Ty_{\Lambda}) \to \pow{\Tm_{\Lambda}(A)^2}\).
  \end{itemize}
  \(\Tm_{\Lambda}(A)\) is the set of terms of type \(A\)
  generated by the language \(\Lambda\), to be defined in \cref{fig:stlc}.

  If \(\Theta = (\Lambda,\mathcal{E})\) we write
  \(\Ty_{\Theta}\) for \(\Ty_{\Lambda}\), and similarly for \(\Tm\).
\end{definition}

The definition of types and terms is given in \cref{fig:stlc}.
\mlb{Technically, we should specify the substitution calculus…}

\begin{figure}
\begin{mathpar}
  \jform{\ctxty{\Gamma}}[Contexts \(\Gamma,\Delta,\dots \in \Ctx_{\Lambda}\)]
  \inferdef[Emp]{ }{\ctxty{\emp}} \and
  \inferdef[Ext]{\ctxty{\Gamma} \\ \ctxty{A}}{\ctxty{\Gamma . (x : A)}} \\
  %
  \jform{\ctxty{A}}[Types \(A,B,\dots \in \Ty_{\Lambda}\)]
  \inferdef[Base]{b \in \mathcal{B}}{\ctxty{b}} \and
  \inferdef[Unit]{ }{\ctxty{\top}} \and
  \inferdef[Prod]{\ctxty{A} \\ \ctxty{B}}{\ctxty{A \times B}} \and
  \inferdef[Fun]{\ctxty{A} \\ \ctxty{B}}{\ctxty{A \to B}} \\
  %
  \jform{\typing{\Gamma}{t}[A]}[Terms \(t,u,\dots \in \Tm_{\Lambda}(A)\)]
  \inferdef[Cst]{c \in \mathcal{C}(A)}{\typing{\Gamma}{c}[A]} \and
  \inferdef[Var]{(x : A) \in \Gamma}{\typing{\Gamma}{x}[A]} \and
  \inferdef[Star]{ }{\typing{\Gamma}{\st}[\top]} \\
  \inferdef[Pair]{\typing{\Gamma}{t}[A] \\ \typing{\Gamma}{u}[B]}{
      \typing{\Gamma}{(t,u)}[A \times B]} \and
  \inferdef[Proj1]{\typing{\Gamma}{p}[A \times B]}{
      \typing{\Gamma}{\proj{1}{p}}[A]} \and
  \inferdef[Proj2]{\typing{\Gamma}{p}[A \times B]}{
      \typing{\Gamma}{\proj{2}{p}}[B]} \\
  \inferdef[Lam]{\typing{\Gamma.(x : A)}{t}[B]}{
      \typing{\Gamma}{\l x.t}[A \to B]} \and
  \inferdef[App]{\typing{\Gamma}{t}[A \to B] \\ \typing{\Gamma}{u}[A]}{
      \typing{\Gamma}{t~u}[B]}
\end{mathpar}
\caption{Types and terms of the simply-typed λ-calculus (with respect to a language
  \(\Lambda = (\mathcal{B},\mathcal{C})\))}
\label{fig:stlc}
\end{figure}

The equational theory is in \cref{fig:stlc-eq}. We omit congruence equations.

\begin{figure}
\begin{mathpar}
  \jform{\conv{\Gamma}{t}{u}[A]}[Equality of terms]
  \inferdef[Refl]{\typing{\Gamma}{t}[A]}{\conv{\Gamma}{t}{t}[A]} \and
  \inferdef[Sym]{\conv{\Gamma}{t}{u}[A]}{\conv{\Gamma}{u}{t}[A]} \and
  \inferdef[Refl]{\conv{\Gamma}{t}{u}[A] \\ \conv{\Gamma}{t}{v}[A]}
    {\conv{\Gamma}{t}{v}[A]} \\
  \text{congruence/substitution rule} \\
  \inferdef[Eq]{(t,u) \in \mathcal{E}(A)}{\conv{\Gamma}{t}{u}[A]} \and
  \inferdef[Unitη]{\typing{\Gamma}{t}[\top]}{\conv{\Gamma}{t}{\st}[\top]} \\
  \inferdef[Pairβ]{\typing{\Gamma}{t}[A] \\ \typing{\Gamma}{u}[B]}{
    \conv{\Gamma}{\proj{1}{(t,u)}}{t}[A] \and
    \conv{\Gamma}{\proj{2}{(t,u)}}{u}[A]
  } \and
  \inferdef[Pairη]{\typing{\Gamma}{p}[A \times B]}{
    \conv{\Gamma}{p}{(\proj{1}{p},\proj{2}{p})}[A \times B]} \\
  \inferdef[Funβ]{\typing{\Gamma,x:A}{t}[B] \\ \typing{\Gamma}{u}[A]}{
    \conv{\Gamma}{(\l x.t)~u}{\subs{t}{x}{u}}[A]} \and
  \inferdef[Funη]{\typing{\Gamma}{f}[A \to B]}{
    \conv{\Gamma}{f}{\l x.(f~x)}[A \to B]}
\end{mathpar}
\caption{Equations of the simply-typed λ-calculus (with respect to a theory
  \(\Theta = (\Lambda,\mathcal{E})\))}
\label{fig:stlc-eq}
\end{figure}

\paragraph{Extra operations}

\begin{definition}{Context partition}{partition}
  The context partition relation is defined as a (partial) relation
  \(\Gamma = \part{\Gamma_{1}}{\Gamma_{2}}\) by
  \begin{mathpar}
    \inferdef{ }{\emp = \part{\emp}{\emp}} \and
    \inferdef{\Gamma = \part{\Gamma_{1}}{\Gamma_{2}}}{
      \Gamma. (x : A) = \part{\Gamma_{1}.(x : A)}{\Gamma_{2}}} \and
    \inferdef{\Gamma = \part{\Gamma_{1}}{\Gamma_{2}}}{
      \Gamma. (x : A) = \part{\Gamma_{1}}{\Gamma_{2}. (x : A)}}
  \end{mathpar}
\end{definition}

\begin{definition}{Reduction}{reduction}

  Reduction is written \(\redop\). \mlb{What sort of reduction is this?}
  
\end{definition}

\subsection{Bidirectional typing}

\begin{figure}
\begin{mathpar}
  \jform{\check{\Gamma}{t}{A}}[Checking/normal terms \(t,u,\dots \in \Nf_{\Lambda}(\Gamma,A)\)]
  
  \inferdef[Star]{ }{\check{\Gamma}{\st}{\top}}
  \label{rule:star-b} \and
  \inferdef[Pair]{\check{\Gamma}{t}{A} \\ \check{\Gamma}{u}{B}}{
      \check{\Gamma}{(t,u)}{A \times B}}
  \label{rule:pair-b} \and
  \inferdef[Lam]{\check{\Gamma,x : A}{t}{B}}{
      \check{\Gamma}{\l x.t}{A \to B}}
  \label{rule:lam-b} \and
  \inferdef[Demote]{\infer{\Gamma}{t}{A} \\ A = B}{\check{\Gamma}{\dem{\icol{t}}}{B}}
  \label{rule:demote-b} \\
  %
  \jform{\infer{\Gamma}{t}{A}}[Infering/neutral terms \(t,u,\dots \in \Ne_{\Lambda}(\Gamma,A)\)]
  \mathcolor{BrickRed}{\inferdef[Cst]{c \in \mathcal{C}(A)}{\infer{\Gamma}{c}{A}}
    } \text{\mlb{Is this infering or checking?}}
  \label{rule:cst-b} \and
  \inferdef[Var]{(x : A) \in \Gamma}{\infer{\Gamma}{x}{A}}
  \label{rule:var-b} \\
  \inferdef[Proj1]{\infer{\Gamma}{p}{A \times B}}{
      \infer{\Gamma}{\pi_1(p)}{A}}
  \label{rule:proj1-b} \and
  \inferdef[Proj2]{\infer{\Gamma}{p}{A \times B}}{
      \infer{\Gamma}{\pi_2(p)}{B}}
  \label{rule:proj2-b} \mprlabel{\textsc{Proj}i} \label{rule:proji-b} \and
  \inferdef[App]{\infer{\Gamma}{t}{A \to B} \\ \check{\Gamma}{u}{A}}{
      \infer{\Gamma}{t~\ccol{u}}{B}}
  \label{rule:app-b}
\end{mathpar}
\caption{Bidirectional typing for the simply-typed λ-calculus
  (with respect to a language \(\Lambda = (\mathcal{B},\mathcal{C})\))}
\label{fig:stlc-bidir}

\end{figure}

We can characterise those terms which have a nice “information flow” by using
bidirectional typing. Incidentally, this characterises normal/neutral forms
– a very deep mystery, but a useful one. Contexts and types are the same as
in \cref{fig:stlc} above, the term judgments are in \cref{fig:stlc-bidir}.

There is an embedding of bidirectional terms into undirected ones, which we
call erasure and write
\begin{align*}
  \eras{\cdot} &\ty \Nf_{\Lambda}(A) \to \Tm_{\Lambda}(A) \\
  \eras{\cdot} &\ty \Ne_{\Lambda}(A) \to \Tm_{\Lambda}(A)
\end{align*}
Both are entirely structural, up to
  \(\eras{\ccol{\dem{\icol{t}}}} = \eras{\icol{t}}\).
We generally leave them implicit.


\begin{theorem}{Normalisation}{thm:norm}
  If \(\Sigma\) is a theory with no equations,
  given a term \(t \in \Tm_{\Sigma}(A)\), there exists
  a unique \(\ccol{n} \in \Nf_{\Sigma}(A)\) such that \(t = \eras{\ccol{n}}\).
\end{theorem}

\subsection{Evaluation context}

\begin{figure}
\begin{mathpar}
  %
  \jform{\stackty{\Gamma}{\pi}{A}{B}}
    [Stacks \(\pi,\dots \in \Sta_{\Lambda}(\Gamma,A,B)\)]
  \inferdef[Nil]{ }{\stackty{\Gamma}{\snil}{T}{T}} \label{rule:nil-s} \and
  \inferdef[Proj\normalfont{i}]{\stackty{\Gamma}{\pi}{A_i}{T}}{
      \stackty{\Gamma}{\sproj{i} \scons \pi}{A_1 \times A_2}{T}}
  \mprlabel{\textsc{Proj}i} \label{rule:proji-s} \and
  \inferdef[App]{\stackty{\Gamma}{\pi}{B}{T} \\ \check{\Gamma}{u}{A}}{
      \stackty{\Gamma}{\sapp{\ccol{u}} \scons \pi}{A \to B}{T}}
  \label{rule:app-s}
\end{mathpar}
\caption{Evaluation stacks for the simply-typed λ-calculus}
\label{fig:stack}

\end{figure}

Another way to capture neutrals, which looks a bit closer to what is happening
in \posscite{Cubric1994} proof, is to use \emph{evaluation stacks},
see \cref{fig:stack}.
A stack represents a succession of elimination forms. Any neutral
can be decomposed into a variable plugged into such a stack, which captures
the succession of eliminations which compose the neutral. The important
difference is that stacks are “inside out”, \ie the top of the stack corresponds
to the innermost elimination of the neutral. This seems to be closer to what
happens in \textcite{Cubric1994}.

\begin{definition}{Stack operations}{stack-op}

Given a stack and a neutral, we can always plug the latter into the former:
\[\begin{array}{rcl}
\plug{\_}{\_} \ty \Sta(\Gamma,A,B) \to \Ne(\Gamma,A) &\to & \Ne(\Gamma,B) \\
\plug{\icol{\snil}}{\icol{t}} &\coloneq & \icol{t} \\
\plug{\icol{(\sproj{i} \scons \pi)}}{\icol{t}} &\coloneq &
  \plug{\icol{\pi}}{\icol{\proj{i}{t}}} \\
\plug{\icol{(\sapp{\ccol{u}} \scons \pi)}}{\icol{t}} &\coloneq &
  \plug{\icol{\pi}}{\icol{t~\ccol{u}}} 
\end{array}\]

Conversely, any neutral can be (uniquely) decomposed into a form
\(\plug{\icol{\pi}}{x}\).

\end{definition}

\section{Interpolation}


\begin{definition}{Polarities}{polarities}
  Let \(A \in \Ty_{\mathcal{B}}\). We define \(A^+ \subseteq \mathcal{B}\)
    (resp. \(A^- \subseteq \mathcal{B}\)) by induction:\pagebreak
  \begin{align*}
    b^{+} & \coloneq \{b\} \\
    b^{-} & \coloneq \emptyset \\
    (A \times B)^{p} & \coloneq A^p \cup B^p \\
    (A \to B)^{p} & \coloneq A^{-p} \cup B^p \\
    \top^p & \coloneq \emptyset
  \end{align*}
  This is lifted to contexts as
  \((\Gamma.(x : A))^p \coloneq \Gamma^p \cup A^p\). 
\end{definition}

\paragraph{Main theorem}

The first version of interpolation works with constant-free theories, \ie those
of the form \(((\mathcal{B},\emptyset),\emptyset)\).

\begin{theorem}{Interpolation – constant-free}{thm:int-inductive}
  Let \(\Sigma\) be a constant-free theory, and assume
  a partitioned context \(\Gamma = \part{\Gamma_{s}}{\Gamma_{t}}\).
  
  Given any normal term
  \(\ccol{t} \in \Nf_{\Sigma}(\Gamma,\tcol{T})\),
  there is a type \(\mcol{M} \in \Ty_{\Sigma}\) and terms
  \(\scol{l} \in \Tm_{\Sigma}(\scol{\Gamma_{s}},\mcol{M})\) and
  \(\tcol{r} \in \Tm_{\Sigma}(\tcol{\Gamma_{t}}.\mcol{(x : M)}, \tcol{T})\)
  such that
  \begin{itemize}
    \item \(\conv{\Gamma}{\subs{\tcol{r}}{x}{\scol{l}}}{\ccol{t}}[\tcol{T}]\);
    \item \(\mcol{M}^+ \subseteq \scol{\Gamma_{s}}^+ \cap (\tcol{\Gamma_{t}}^- \cup \tcol{T}^+)\) and \(\mcol{M}^- \subseteq \scol{\Gamma_{s}}^- \cap (\tcol{\Gamma_{t}}^+ \cup \tcol{T}^-)\).
  \end{itemize}
  
  Moreover, given any neutral term
  \(\icol{t} \in \Ne_{\Sigma}(\Gamma,T)\), there is a type
  \(\mcol{M} \in \Ty_{\Sigma}\) such that either:
  \begin{enumerate}
    \item “\(\scol{T}\) is a source type”, \ie
      we have terms \(\tcol{l} \in \Tm_{\Sigma}(\tcol{\Gamma_{t}},\mcol{M})\) and
  \(\scol{r} \in \Tm_{\Sigma}(\scol{\Gamma_{s}}.\mcol{(x : M)}, \scol{T})\) such that
      \begin{itemize}
        \item \(\conv{\Gamma}{\subs{\scol{r}}{x}{\tcol{l}}}{\icol{t}}[\scol{T}]\),
        \item \(\scol{T}^+ \subseteq \scol{\Gamma_{s}}^+\) and
          \(\scol{T}^- \subseteq \scol{\Gamma_{s}}^-\),
        \item \(\mcol{M}^+ \subseteq \scol{\Gamma_{s}}^+ \cap (\tcol{\Gamma_{t}}^- \cup \tcol{T}^+)\)
          and \(\mcol{M}^- \subseteq \scol{\Gamma_{s}}^- \cap (\tcol{\Gamma_{t}}^+ \cup \tcol{T}^-)\);
      \end{itemize}
    \item or “\(\tcol{T}\) is a target type”, \ie
      we have terms \(\scol{l} \in \Tm_{\Sigma}(\scol{\Gamma_{s}},\mcol{M})\) and
      \(\tcol{r} \in \Tm_{\Sigma}(\tcol{\Gamma_{t}}.\mcol{(x : M)}, \tcol{T})\) such that
      \begin{itemize}
        \item \(\conv{\Gamma}{\subs{\tcol{r}}{x}{\scol{l}}}{\icol{t}}[\scol{T}]\),
        \item \(\tcol{T}^+ \subseteq \tcol{\Gamma_{t}}^+\) and
          \(\tcol{T}^- \subseteq \tcol{\Gamma_{t}}^-\),
        \item \(\mcol{M}^+ \subseteq \scol{\Gamma_{s}}^+ \cap (\tcol{\Gamma_{t}}^- \cup \tcol{T}^+)\)
          and \(\mcol{M}^- \subseteq \scol{\Gamma_{s}}^- \cap (\tcol{\Gamma_{t}}^+ \cup \tcol{T}^-)\).
      \end{itemize}
  \end{enumerate}
\end{theorem}

\begin{proof}
  By induction on the definition of \(\Ne/\Nf\).
  
  The \(\Nf\) cases are rather direct.

  \nameref{rule:star-b}: \(\tcol{T}\) is \(\top\), so we can take
  \(\mcol{M} = \top\), and \(\scol{l} = \tcol{r} = \star\).

  \nameref{rule:pair-b}: \(\tcol{T}\) is \(\tcol{T_1} \times \tcol{T_2}\), \(\ccol{t}\)
  is \(\ccol{(t_1,t_2)}\). By induction,
  we obtain \(\scol{l_i} \in \Tm(\scol{\Gamma_{s}},\mcol{M_i})\)
  and \(\scol{r_i} \in \Tm(\tcol{\Gamma_{t}}.\mcol{(x : M_i)},\tcol{T_i})\).
  We can then take \(\mcol{M} \coloneq \mcol{M_1 \times M_2}\) and
  \[\begin{array}{rlll}
    \scol{\Gamma_{s}} \vdash & \scol{l} \coloneq &
      \scol{(l_1,l_2)} & \ty \mcol{M_1 \times M_2} \\
    \tcol{\Gamma_{t}}.\mcol{(x : M_1 \times M_2)} \vdash &
      \tcol{r} \coloneq & \tcol{(\subs{r_1}{y}{\proj{1}{y}},\subs{r_2}{y}{\proj{2}{y}})} &
      \ty \tcol{T_1 \times T_2}
  \end{array}\]
  
  \nameref{rule:lam-b}: We extend the \emph{target} context with the extra variable
  (this is why we need this more complicated formulation with a context split!)
  and by induction hypothesis, we obtain
  \(\scol{l'} \in \Tm(\scol{\Gamma_{s}},\mcol{M})\) and
  \(\tcol{r'} \in \Tm(\tcol{\Gamma_{t}}.\mcol{(x : M)}.\tcol{(y : A)})\),
  and we can take
  \[\begin{array}{rlll}
  \scol{\Gamma_{s}} \vdash & \scol{l} \coloneq & \scol{l'} & \ty \mcol{M} \\
  \tcol{\Gamma_{t}}.\mcol{(x : M)} \vdash & \tcol{r} \coloneq &
    \tcol{\l y.r'} &
    \ty \tcol{A \to B}
  \end{array}\]


  The \(\Ne\) cases are more interesting. For each rule, we have two subcases,
  depending on the source/target induction hypothesis for the main sub-neutral,
  or on whether the variable in is the source or target context.

  \nameref{rule:var-b}/1: We have \((x \ty \scol{A}) \in \scol{\Gamma_s}\).
  Then we take \(\mcol{M} \coloneq \top\)
  \[\begin{array}{rlll}
  \tcol{\Gamma_{t}} \vdash & \tcol{l} \coloneq & \tcol{\star} & \ty \mcol{M} \\
  \scol{\Gamma_{s}}.\mcol{(\_ : M)} \vdash & \scol{r} \coloneq &
    \scol{x} & \ty \scol{A}
  \end{array}\]

  \nameref{rule:var-b}/2: We have \((x \ty \tcol{A}) \in \tcol{\Gamma_t}\).
  Then we take \(\mcol{M} \coloneq \top\)
  \[\begin{array}{rlll}
  \scol{\Gamma_{s}} \vdash & \scol{l} \coloneq & \scol{\star} & \ty \mcol{M} \\
  \tcol{\Gamma_{t}}.\mcol{(\_ : M)} \vdash & \tcol{r} \coloneq &
    \tcol{x} & \ty \tcol{A}
  \end{array}\]

  \nameref{rule:demote-b}/1: By induction hypothesis, we get that \(\scol{A}\)
    is a “source type”,
    \(\tcol{l} \in \Tm_{\Sigma}(\tcol{\Gamma_{t}},\mcol{M})\) and
    \(\scol{r} \in \Tm_{\Sigma}(\scol{\Gamma_{s}}.\mcol{(x : M)}, \scol{A})\).
    But because of the constraint of the \nameref{rule:demote-b} rule, it must
    be the case that \(\scol{A} = \tcol{B}\), in other words these type really
    are in the common language of \(\scol{\Gamma_s}\) and \(\tcol{\Gamma_t}\).
    Thus, we take \(\mcol{M'} \coloneq \mcol{M \to A}\) and
    \[\begin{array}{rlll}
      \scol{\Gamma_{s}} \vdash & \scol{l'} \coloneq & \l x.\scol{r} & \ty \mcol{M \to A} \\
      \tcol{\Gamma_{t}}.\mcol{(z : M)} \vdash & \tcol{r'} \coloneq &
        \tcol{z~l} & \ty \tcol{A}
    \end{array}\]
  \nameref{rule:demote-b}/2: In this case we keep the exact
    same interpolating type and interpolants.
    
  \nameref{rule:app-b}/1: We are in the case of an application
  \(\icol{t}~\ccol{u}\), with the induction hypothesis on \(\icol{t}\) giving
  us some \(\scol{A \to B}\) with vocabulary in \(\scol{\Gamma_s}\) and
  \(\mcol{M_t}\), \(\scol{r_t}\) and \(\tcol{l_t}\) such that
  \(\icol{t} = \subs{\scol{r_t}}{x}{\tcol{l_t}}\). Because \(\scol{A}\) has its
  vocabulary in \(\scol{\Gamma_s}\), by induction hypothesis we can
  interpolate \(\ccol{u}\) with a “reversed” colouring on \(\Gamma\),
  \ie inverting the roles of \(\scol{\Gamma_s}\) and \(\tcol{\Gamma_t}\).
  This gives us \(\mcol{M_u}\), \(\scol{r_u}\) and \(\tcol{l_u}\) such that
  \(\ccol{u} = \subs{\scol{r_u}}{x}{\tcol{l_u}}\). Then we take
  \(\mcol{M} \coloneq \mcol{M_t} \times \mcol{M_u}\), and
  \[\begin{array}{rlll}
      \tcol{\Gamma_{t}} \vdash & \tcol{l} \coloneq & \tcol{(l_t,l_u)} &
        \ty \mcol{M_t \times M_u} \\
      \scol{\Gamma_{t}}.\mcol{(z : M_t \times M_u)} \vdash & \scol{r} \coloneq &
        \scol{\subs{r_t}{x}{\proj{1}{z}}~\subs{r_u}{x}{\proj{2}{z}}} & \ty \scol{B}
  \end{array}\]

  \nameref{rule:app-b}/2: This time, the induction hypothesis on \(\icol{t}\) gives
  us some \(\tcol{A \to B}\) with vocabulary in \(\tcol{\Gamma_s}\) and
  \(\mcol{M_t}\), \(\tcol{r_t}\) and \(\scol{l_t}\) such that
  \(\icol{t} = \subs{\tcol{r_t}}{x}{\scol{l_t}}\). We can directly use the
  induction hypothesis on \(\ccol{u}\), to get
  \(\mcol{M_u}\), \(\tcol{r_u}\) and \(\scol{l_u}\) such that
  \(\ccol{u} = \subs{\tcol{r_u}}{x}{\scol{l_u}}\). Then we take again
  \(\mcol{M} \coloneq \mcol{M_t} \times \mcol{M_u}\), and
  \[\begin{array}{rlll}
      \scol{\Gamma_{s}} \vdash & \scol{l} \coloneq & \scol{(l_t,l_u)} &
        \ty \mcol{M_t \times M_u} \\
      \tcol{\Gamma_{t}}.\mcol{(z : M_t \times M_u)} \vdash & \tcol{r} \coloneq &
        \tcol{\subs{r_t}{x}{\proj{1}{z}}~\subs{r_u}{x}{\proj{2}{z}}} & \ty \tcol{B}
  \end{array}\]

  \nameref{rule:proji-b}: Here in all cases we keep the same interpolating type
    and left interpolant, and simply post-compose the right interpolant with
    the relevant projection.
    
\end{proof}

It seems we do not quite construct the same interpolant as \citeauthor{Cubric1994}.
Indeed, in the case of a stack of applications on a source variable of type
\(\scol{A_1 \to \dots \to A_n \to B}\), we will construct an interpolating type
\(\mcol{(M_1 \times \dots \times M_n) \to B}\), while \citeauthor{Cubric1994} will
rather construct \(\mcol{M_1} \to \dots \to \mcol{M_n} \to \mcol{B}\).

\printbibliography

\end{document}